\chapter{Lumbago}
\index{Lumbago}

\section*{Definition and Overview}
Lumbago is a broad clinical term used to describe acute or chronic pain located in the lumbar region of the back, commonly referred to as lower back pain. It is not a specific disease entity but rather a symptom arising from a wide array of underlying musculoskeletal, neurological, or systemic conditions. The pain can vary from a dull, constant ache to a sudden, sharp, debilitating sensation, often accompanied by stiffness and limited mobility. Given its high prevalence and potential for causing substantial disability and occupational impairment, lumbago is a major public health concern with significant socioeconomic implications.

\section*{Epidemiology}
Lower back pain is one of the most common ailments globally, with a lifetime prevalence estimated at approximately 80\% to 85\% in the general population. It is the leading cause of years lived with disability (YLDs) worldwide. In many countries, millions of people suffer from lower back pain, affecting individuals across all age groups but peaking in prevalence between 35 and 55 years. Musculoskeletal lumbago is equally common in males and females, although certain underlying aetiologies like osteoporosis-related fractures are more common in elderly females. The direct and indirect economic costs of lumbago, including healthcare utilisation and lost productivity, represent a multi-billion-dollar burden annually.

\section*{Aetiology and Pathophysiology}
Lumbago can be caused by mechanical disorders of the spinal column, neurological impingement, or, less frequently, serious systemic disease.
\begin{itemize}
    \item \textbf{Mechanical strain:} The most common cause, involving acute micro-tears or strain of the paraspinal muscles, tendons, or ligaments of the lumbar spine, typically triggered by lifting heavy objects, sudden twisting, or poor posture.
    \item \textbf{Lumbar disc pathology:} Herniation or degeneration of the intervertebral discs (degenerative disc disease), which reduces shock absorption, alters spinal mechanics, and can directly irritate nearby nociceptors.
    \item \textbf{Radiculopathy (Sciatica):} Compression or inflammation of a spinal nerve root, often by a herniated disc or osteophyte, leading to radiating pain along the path of the sciatic nerve.
    \item \textbf{Facet joint arthropathy:} Osteoarthritis of the lumbar facet joints, leading to localised inflammation, stiffness, and pain that worsens with extension and twisting of the spine.
    \item \textbf{Spinal stenosis:} Age-related narrowing of the spinal canal or neural foramina, causing compression of the cauda equina or nerve roots, classically presenting as neurogenic claudication.
    \item \textbf{Spondylolisthesis:} The forward slippage of one vertebra over another, which disrupts spinal alignment and can stretch nerve roots or compress the spinal canal.
    \item \textbf{Systemic and referred causes:} Serious non-mechanical causes such as vertebral osteomyelitis, epidural abscess, metastatic malignancy, inflammatory spondyloarthropathies (e.g., ankylosing spondylitis), or referred pain from abdominal aortic aneurysm or nephrolithiasis.
\end{itemize}

\section*{Clinical Features}
The presentation of lumbago depends on the underlying cause, but typically involves pain and mechanical limitation in the lower back region.
\begin{itemize}
    \item \textbf{Localised lumbar pain:} Dull, aching, or sharp pain confined to the lower back, which may radiate to the buttocks or upper thighs but does not extend below the knee in purely mechanical cases.
    \item \textbf{Radicular pain (Sciatica):} Sharp, shooting, or electric-shock-like pain that radiates from the back down the back or side of one leg, extending below the knee, often accompanied by numbness, tingling (paresthesia), or focal weakness.
    \item \textbf{Spinal stiffness and muscle spasm:} Significant restriction in spinal range of motion, particularly during flexion, extension, or rotation, accompanied by palpable, painful spasms of the paraspinal muscles.
    \item \textbf{Neurogenic claudication:} Lower limb pain, numbness, or weakness that is brought on by walking or standing and relieved by sitting or bending forward (the "shopping trolley sign"), characteristic of spinal stenosis.
    \item \textbf{Postural abnormalities:} Antalgic gait or leaning to one side to minimise pain, and difficulty transitioning from sitting to standing.
    \item \textbf{Red flag symptoms:} Severe night pain, unexplained weight loss, fever, history of malignancy, or sudden bowel or bladder dysfunction, which indicate a non-mechanical or emergency aetiology.
\end{itemize}

\section*{Differential Diagnosis}
Due to the non-specific nature of lower back pain, it is essential to differentiate benign mechanical lumbago from serious spinal or systemic diseases.
\begin{itemize}
    \item \textbf{Cauda equina syndrome:} A surgical emergency caused by massive compression of the lumbosacral nerve roots, presenting with saddle anaesthesia, bladder/bowel incontinence or retention, and bilateral leg weakness, unlike typical lumbago.
    \item \textbf{Ankylosing spondylitis:} An inflammatory arthritis presenting with chronic lower back pain that is worse in the morning, improves with exercise, and is associated with sacroiliitis on imaging.
    \item \textbf{Renal colic (Nephrolithiasis):} Severe, colicky loin-to-groin pain, often accompanied by haematuria, nausea, and vomiting, which can mimic acute lumbago but has no mechanical spinal triggers.
    \item \textbf{Abdominal aortic aneurysm (AAA) rupture:} Presents with sudden, tearing back or abdominal pain, hypotension, and a pulsatile abdominal mass, requiring immediate emergency intervention.
    \item \textbf{Spinal metastasis:} Pain that is constant, progressive, unrelieved by rest, and occurs in a patient with a history of cancer.
\end{itemize}

\section*{When to Seek Medical Care}
Most cases of mechanical lumbago improve with self-care, but prompt medical evaluation is necessary if the pain is severe, persistent, or accompanied by neurological deficits.

\textbf{Contact your local emergency services for:}
\begin{itemize}
    \item Sudden loss of bowel or bladder control (incontinence or acute retention).
    \item Numbness in the groin or saddle area (saddle anaesthesia) or sudden weakness in both legs making it impossible to walk.
    \item Sudden, severe, tearing back or abdominal pain accompanied by dizziness, fainting, or collapse.
\end{itemize}
Other non-emergency reasons to see a doctor include pain that does not improve after 4 to 6 weeks, pain accompanied by fever or unexplained weight loss, or radicular pain radiating below the knee.

\section*{Diagnosis and Investigations}
The diagnosis of lumbago relies on a detailed history and physical examination, with routine imaging discouraged for acute, uncomplicated back pain unless "red flags" are present.
\begin{itemize}
    \item \textbf{Clinical history and physical examination:} Evaluating pain characteristics, looking for red flags, and performing neurological testing (reflexes, sensation, motor strength) and a straight leg raise test (to screen for nerve root irritation).
    \item \textbf{Lumbar spine X-ray:} Helpful only to evaluate for fractures (especially in patients with osteoporosis or trauma) or gross structural instability, but does not show soft tissue or disc pathology.
    \item \textbf{Magnetic resonance imaging (MRI):} The imaging modality of choice for suspected radiculopathy, spinal stenosis, infection, or malignancy, showing detailed views of intervertebral discs, nerve roots, and soft tissues.
    \item \textbf{Computed tomography (CT) scan:} Used as an alternative to MRI when MRI is contraindicated or to assess bony anatomy in detail.
    \item \textbf{Inflammatory markers (CRP and ESR):} Ordered when inflammatory arthritis, spinal infection, or malignancy is suspected based on clinical features.
\end{itemize}

\section*{Management}
The management of lumbago focuses on maintaining activity, providing symptomatic relief, and avoiding prolonged bed rest.
\begin{itemize}
    \item \textbf{Non-pharmacological measures:} Staying active and continuing daily activities as tolerated, applying heat packs to reduce muscle spasm, and avoiding bed rest for more than 24 to 48 hours.
    \item \textbf{Pharmacological analgesia:} First-line treatment with non-steroidal anti-inflammatory drugs (NSAIDs) like ibuprofen or naproxen to reduce pain and inflammation, with paracetamol or short-term muscle relaxants as adjunctive options.
    \item \textbf{Physiotherapy and exercise:} Active rehabilitation including stretching, core-strengthening exercises, and aerobic conditioning once the acute pain subsides, which reduces the risk of chronic pain.
    \item \textbf{Epidural steroid injections:} Targeted injection of corticosteroids and local anaesthetic into the epidural space, reserved for patients with severe, refractory radicular pain (sciatica) to reduce nerve root inflammation.
    \item \textbf{Surgical intervention:} Decompressive surgeries (such as laminectomy or discectomy) or spinal fusion, reserved for patients with progressive neurological deficits, cauda equina syndrome, or debilitating pain that fails to respond to extensive conservative therapy.
    \item \textbf{Psychosocial support:} Cognitive behavioural therapy (CBT) and addressing "yellow flags" (beliefs, behaviours, and social factors that predict chronic disability) in patients with long-standing, chronic lumbago.
\end{itemize}

\section*{Complications}
\begin{itemize}
    \item \textbf{Chronic pain syndrome:} Persistent lower back pain lasting beyond 12 weeks, which can cause significant functional limitations and psychological distress.
    \item \textbf{Progressive motor weakness:} Permanent nerve damage leading to foot drop or muscle wasting in the lower limbs if nerve root compression is left untreated.
    \item \textbf{Cauda equina damage:} Permanent bladder, bowel, or sexual dysfunction if severe compression of the lumbosacral nerve roots is not surgically decompressed in a timely manner.
    \item \textbf{Analgesic adverse effects:} Gastric ulceration, renal impairment, or medication dependence issues arising from the inappropriate long-term use of pain medications.
    \item \textbf{Loss of employment and disability:} Significant occupational disruption, loss of income, and reduced social engagement due to chronic physical limitations.
\end{itemize}

\section*{Prognosis and Outcomes}
The prognosis for acute mechanical lumbago is highly favourable, with approximately 90\% of patients experiencing substantial recovery within 4 to 6 weeks, regardless of the treatment modality used. However, recurrence is common, with up to 50\% of patients experiencing another episode within a year. A small proportion (5\% to 10\%) go on to develop chronic lower back pain, which is more challenging to treat and is strongly influenced by occupational and psychosocial factors rather than radiological severity.

\section*{Prevention and Self-Care}
\begin{itemize}
    \item \textbf{Maintain a healthy weight:} Reducing excess abdominal fat to decrease mechanical strain on the lumbar spine.
    \item \textbf{Regular core exercises:} Engaging in activities like Pilates, yoga, or swimming to strengthen abdominal and back muscles, which support the spine.
    \item \textbf{Proper lifting techniques:} Bend the knees and lift with the legs rather than the back, keeping the load close to the body, and avoid twisting while lifting.
    \item \textbf{Ergonomic workstation setup:} Ensure chairs provide adequate lumbar support, and take frequent breaks to stand and stretch if sitting for long periods.
    \item \textbf{Quit smoking:} Tobacco use reduces blood flow to the intervertebral discs, accelerating degenerative changes and increasing back pain risk.
    \item \textbf{Choose supportive footwear:} Avoid high heels and wear flat, comfortable shoes to maintain proper pelvic alignment.
\end{itemize}

\section*{Sources and Further Reading}
The following reputable organisations provide further information and clinical guidelines on lower back pain and lumbago:
\begin{itemize}
    \item Healthdirect Australia. \textit{Lower back pain symptoms, causes, and treatments.} https://www.healthdirect.gov.au/
    \item Better Health Channel. \textit{Back pain - causes, treatment, and prevention.} https://www.betterhealth.vic.gov.au/
    \item National Institute of Arthritis and Musculoskeletal and Skin Diseases (NIAMS). \textit{Low back pain information page.} https://www.niams.nih.gov/
    \item National Institute for Health and Care Excellence (NICE). \textit{Low back pain and sciatica in over 16s: assessment and management.} https://www.nice.org.uk/
    \item Mayo Clinic. \textit{Back pain causes, diagnosis, and treatment.} https://www.mayoclinic.org/
\end{itemize}
